\subsection{Training Pipeline}
In this part, the training pipeline of different 3D resources are introduced.

To specify, the whole experimental environment is built upon a single NVIDIA GeForce RTX 4090 GPU, which features 24 GB memory. The processing unit comprises 20 virtual CPU cores (vCPUs) allocated from an Intel Xeon Platinum 8470Q server processor, while the system is supplemented with 90 GB of available random-access memory to accommodate large-scale data loading and computational overheads during training and inference from \textbf{AutoDL}
\subsubsection{Object A}
The real-world multi-view object is implemented based on the 2DGS framework, which consists of 2 major stages: camera pose estimation, and Gaussian Splatting optimization.

First, COLMAP is used to perform structure-from-motion (SfM) reconstruction and estimate camera intrinsics and extrinsics. Sparse point clouds and camera poses generated by COLMAP are then used to initialize the Gaussian representation.

Then, the reconstructed scene is optimized using the official Gaussian Splatting framework, which generates the original point cloud.



\subsubsection{Object B}
The text-driven 3D asset generation task is implemented based on the \textit{threestudio} framework, which adopts the newly-released T2I model DeepFloyd IF and SDS loss. Given a textual prompt, the framework initializes a learnable 3D representation and iteratively optimizes it by minimizing the discrepancy between rendered views and diffusion-guided image priors.

Specifically, a pretrained text-to-image diffusion model provides semantic guidance during optimization. At each iteration, SDS loss is computed to encourage consistency with the textual description.

\subsubsection{Object C}
The single-image 3D generation task is implemented using the Magic123 framework, which consists of coarse geometry initialization, diffusion-guided optimization.

First,a Neural Radiance Field (NeRF) is optimized from scratch using the input RGBA image and a text prompt describing the object.

Then,the coarse NeRF is first converted into a textured mesh via a differentiable marching tetrahedra (DMTet) representation, using a density threshold to extract the surface.
Starting from this initialization, the mesh geometry and texture are further refined with the same Zero123 and Stable Diffusion guidance.